\documentclass[10pt,letterpaper]{article}
% Shared corporate style for StrategicDefense Systems Security LTD.
\usepackage[T1]{fontenc}
\usepackage[utf8]{inputenc}
\usepackage{tgtermes}
\usepackage{microtype}
\usepackage{xcolor}
\usepackage{graphicx}
\usepackage{geometry}
\usepackage{fancyhdr}
\usepackage{lastpage}
\usepackage{enumitem}
\usepackage{tabularx}
\usepackage{array}
\usepackage{booktabs}
\usepackage{setspace}
\usepackage{ragged2e}
\usepackage[hidelinks]{hyperref}

\definecolor{SDSSNavy}{HTML}{173C64}
\definecolor{SDSSGreen}{HTML}{183C31}
\definecolor{SDSSRule}{HTML}{A9B8C8}

\newcommand{\CompanyFull}{StrategicDefense Systems Security LTD.}
\newcommand{\CompanyShort}{Strategic Defense}
\newcommand{\ContactLine}{4 Oguda Street, AMAC, Abuja, FCT, Nigeria\enspace\textbar\enspace Email: info@strategic-defence.com\enspace\textbar\enspace Telephone: +234 703 867 6619}
\graphicspath{{../../}}

\hypersetup{
  colorlinks=true,
  linkcolor=SDSSNavy,
  urlcolor=SDSSNavy,
  pdfauthor={StrategicDefense Systems Security LTD.}
}

\newcommand{\SDSSLogo}[1][2.6in]{\includegraphics[width=#1]{logo_transparent.png}}
\newcommand{\SDSSIcon}[1][0.15in]{\includegraphics[width=#1]{icon 1.png}}

\geometry{left=0.72in,right=0.72in,top=0.48in,bottom=0.48in}
\setlength{\parindent}{0pt}
\setlength{\parskip}{4.3pt}
\setstretch{0.96}
\pagestyle{empty}

\hypersetup{
  pdftitle={Submission of Preliminary Technical Proposal - SHIELD},
  pdfsubject={Secure Homeland Integrated Early-Response and Land-Border Defence}
}

\begin{document}
\begin{center}
  \SDSSLogo[2.75in]\\[6pt]
  {\small \ContactLine}
\end{center}

\vspace{2pt}
{\color{SDSSNavy}\rule{\textwidth}{0.8pt}}
\vspace{5pt}

\begin{tabular*}{\textwidth}{@{\extracolsep{\fill}}lll}
  Our Ref: SDSS/SHIELD/CL/001 & Version: 1.0 & Date: 26 September 2026
\end{tabular*}

\vspace{7pt}
Major-General Adeyinka A. Famadewa (Rtd.)\\
Special Adviser to the President on Homeland Security\\
Office of the Special Adviser to the President on Homeland Security\\
Abuja, Nigeria

\vspace{7pt}
{\bfseries SUBJECT: SUBMISSION OF PRELIMINARY TECHNICAL PROPOSAL - SHIELD (SECURE HOMELAND INTEGRATED EARLY-RESPONSE AND LAND-BORDER DEFENCE)}

\vspace{4pt}
Dear Sir,

\CompanyShort{} submits its Preliminary Technical Proposal for SHIELD, an equipment-based capability to strengthen border management and joint emergency response in support of the homeland security priorities of your office.

The proposal builds on the coordination already established through the Intelligence Fusion Centre and extends it to where it is needed most: the land borders and the scene of an emergency. It reflects the principle that security is not the responsibility of the military alone, but is shared by civilian and military agencies and by communities. We offer two capability packages. Package A gives customs, immigration, police and civil defence a continuous view of a pilot section of the South-West border with Benin, to curb smuggling and trafficking, protect border communities and speed up legitimate trade. Package B provides interoperable communications and a mobile incident command vehicle, so that police, emergency services, civil defence and, where needed, the military can talk and coordinate as one team at any major incident, from attacks to floods and building collapses, without replacing their own radios.

The proposal rests on a single principle: shared picture, not shared command. Each agency retains its mandate and chain of command throughout. SHIELD equipment is technology-neutral and built on open standards, so that it connects to national systems and avoids dependence on a single supplier.

This is a preliminary proposal, intended as a basis for technical discussion with your office. The Final Technical Proposal will follow a site survey and will confirm quantities, locations and schedule. Delivery would begin with an early live demonstration of cross-agency radio interoperability.

\CompanyShort{} is ready to present the proposal to your office at a time convenient to you, and to agree the pilot border section and participating agencies as the next step.

\vspace{2pt}
Yours faithfully,

\vspace{16pt}
\rule{2.35in}{0.4pt}\\[-1pt]
{\color{SDSSNavy}\bfseries Wael ALTAKI}\\[-1pt]
General Manager\\[-1pt]
\CompanyFull

\vspace{5pt}
{\color{SDSSNavy}\bfseries ENCLOSURES}

\begin{enumerate}[leftmargin=1.6em,itemsep=0pt,topsep=1pt]
  \item Preliminary Technical Proposal, Ref: SDSS/SHIELD/PTP/001 v1.0
\end{enumerate}
\end{document}
